% Minimal preamble for Open Logic Project sources (pinned commit 1e960beff9ed7835bf3e3f1335e21af3439cd107).
%
% OLP's own master preamble cannot run on the project host: TeX Live 2025 there lacks
% prooftree.sty, ntheorem.sty, extarrows.sty, mathdots.sty, isomath.sty and subfiles.sty.
% Following the project's own rule ("if the author preamble needs packages the toolchain lacks,
% do not invent macros -- use a minimal working preamble"), this file declares, verbatim in
% spirit and numbering behaviour, exactly the theorem-like environments OLP defines in
% sty/open-logic-book-envs.sty (thm/lem/prop/cor/defn/ex/prob/rem/note/case/axiom/conv),
% plus the shorthand accents OLP uses (`!!a`, `!!b`, `!!{...}`).
%
% Author text is untouched; only the packaging is ours.
\documentclass[10pt,a4paper]{article}

\usepackage{amsmath}
\usepackage{amssymb}
\usepackage{amsthm}
\usepackage{graphicx}
\usepackage{hyperref}
\usepackage{enumitem}
\usepackage[framemethod=default]{mdframed}

% ---- OLP shorthand accents (from sty/open-logic.sty) ----
% NB: \! is taken by LaTeX (negative thin space), so we only define the !!a / !!b families
% and the \olup helper used by them.
\newcommand*{\olup}[1]{\MakeUppercase{#1}}
\def\!!!a#1{\olup{#1}}
\def\!a#1{\olup{#1}}
\def\!b#1{\olup{#1}}

% ---- theorem-like environments, same family numbering as OLP ----
\theoremstyle{plain}
\newtheorem{thm}{Theorem}[section]
\newtheorem{lem}[thm]{Lemma}
\newtheorem{prop}[thm]{Proposition}
\newtheorem{cor}[thm]{Corollary}
\newtheorem{axiom}{Axiom}

\theoremstyle{definition}
\newtheorem{defn}[thm]{Definition}
\newtheorem{ex}[thm]{Example}
\newtheorem{prob}{Problem}[section]
\newtheorem{probd}{Problem}
\newtheorem{conv}[thm]{Convention}

\theoremstyle{remark}
\newtheorem{rem}[thm]{Remark}
\newtheorem{note}[thm]{Note}
\newtheorem{case}[thm]{Case}

% ---- environments OLP defines as wrappers ----
\newenvironment{probtag}[1]{\begin{prob}}{\end{prob}}
\newenvironment{explain}{}{}
\newenvironment{explain*}{}{}
\newenvironment{editorial}{}{}
\newenvironment{digress}{}{}
\newenvironment{digress*}{}{}
\newenvironment{history}{}{}
\newenvironment{reference}{}{}
\newcommand*{\olfileid}[3]{}
\newcommand*{\olsection}[1]{\section*{#1}}
\newcommand*{\olchapter}[3]{}
\newcommand*{\olpart}[2]{}
\newcommand*{\olimport}[1]{}
\newcommand*{\olinclude}[1]{}
\newcommand*{\olpath}[1]{}
\newcommand*{\ollangid}[1]{}
\newcommand*{\problemtag}[1]{}
\newcommand*{\problemauthor}[1]{}
\newcommand*{\ollabel}[1]{\label{#1}}
% \olref takes 0-3 OPTIONAL arguments plus one mandatory label:
%     \olref{lab}  |  \olref[from]{lab}  |  \olref[from][to]{lab}  |  \olref[a][b][c]{lab}
% A single-argument definition silently swallows "[sth]" and emits a broken reference
% (measured on the host: the text came out as "??iso]x"). The chain below consumes the
% brackets and cross-references only the mandatory argument.
\makeatletter
\newcommand{\olref}{\@ifnextchar[{\olref@i}{\olref@final}}
\def\olref@i[#1]{\@ifnextchar[{\olref@ii}{\olref@final}}
\def\olref@ii[#1]{\@ifnextchar[{\olref@iii}{\olref@final}}
\def\olref@iii[#1]{\olref@final}
\def\olref@final#1{\ref{#1}}
\makeatother
\newcommand*{\olnewpage}{\clearpage}

% ===================== OLP compatibility layer (1) =====================

\makeatletter

% ---------------- proof trees (OLP's prooftree is a bussproofs derivative) ----------------
\providecommand{\AxiomC}[1]{}
\providecommand{\UnaryInfC}[1]{}
\providecommand{\BinaryInfC}[1]{}
\providecommand{\TrinaryInfC}[1]{}
\providecommand{\QuaternaryInfC}[1]{}
\providecommand{\QuinaryInfC}[1]{}
\providecommand{\RightLabel}[1]{}
\providecommand{\LeftLabel}[1]{}
\providecommand{\doubleLine}{}
\providecommand{\noLine}{}
\providecommand{\DisplayProof}{}
\providecommand{\fCenter}{}
\newenvironment{prooftree}{}{}
\newenvironment{ltprooftree}{}{}
\newenvironment{oltableau}{}{}

% ---------------- logical notation ----------------
\providecommand{\lif}{\mathbin{\rightarrow}}
\providecommand{\lforall}{\mathord{\forall}}
\providecommand{\lexists}{\mathord{\exists}}
\providecommand{\ltrue}{\top}
\providecommand{\lfalse}{\bot}
\providecommand{\Ltrue}{\top}
\providecommand{\Lfalse}{\bot}
\providecommand{\Sequent}{\mathrel{\Rightarrow}}
\providecommand{\Proves}{\mathrel{\vdash}}
\providecommand{\Entails}{\mathrel{\models}}
\providecommand{\Subst}[3]{#1[#2/#3]}
\providecommand{\sFmla}[2]{#1(#2)}
\providecommand{\pFmla}[2]{#1(#2)}
\providecommand{\formula}{\relax}
\providecommand{\True}{\top}
\providecommand{\False}{\bot}
\providecommand{\Args}[1]{#1}
\providecommand{\Ts}{\ensuremath{\top}}
\providecommand{\Fi}{\ensuremath{\bot}}
\providecommand{\lfalsum}{\bot}

% ---------------- objects, structures, metatheory ----------------
\providecommand{\Obj}{}
\providecommand{\Struct}{\ensuremath{\mathfrak{A}}}
\providecommand{\Domain}{\ensuremath{|\mathfrak{A}|}}
\providecommand{\Assign}{\ensuremath{s}}
\providecommand{\Sat}[2]{\ensuremath{#1 \models #2}}
\providecommand{\mSat}[3]{\ensuremath{#1 \models #2[#3]}}
\providecommand{\mModel}[2]{\ensuremath{#1 \models #2}}
\providecommand{\Lang}{\ensuremath{\mathcal{L}}}
\providecommand{\Th}{\ensuremath{\mathrm{Th}}}
\providecommand{\Log}{\ensuremath{\mathrm{L}}}
\providecommand{\tuple}[1]{\langle #1 \rangle}
\providecommand{\num}[1]{\underline{#1}}
\providecommand{\cfind}[1]{\varphi_{#1}}
\providecommand{\Nat}{\ensuremath{\mathbb{N}}}
\providecommand{\Mult}{\ensuremath{\mathrm{Mult}}}
\providecommand{\Add}{\ensuremath{\mathrm{Add}}}
\providecommand{\Exp}{\ensuremath{\mathrm{Exp}}}
\providecommand{\Zero}{\ensuremath{0}}
\providecommand{\Succ}{\ensuremath{\mathrm{S}}}
\providecommand{\lambd}{\ensuremath{\lambda}}
\providecommand{\bredpar}{\mathrel{\rightarrow_\beta}}
\providecommand{\bcd}[1]{\ensuremath{\ulcorner #1 \urcorner}}
\providecommand{\Set}[1]{\ensuremath{\{#1\}}}
\providecommand{\indcase}[1]{\par\noindent\emph{Case #1}.\quad}
\providecommand{\indfrm}[1]{\par\noindent\emph{Form #1}.\quad}
\providecommand{\iftag}[1]{#1}
\providecommand{\tagitem}[1]{\item[#1]}
\providecommand{\Elim}[1]{}
\providecommand{\Intro}[1]{}
\providecommand{\LeftR}[1]{}
\providecommand{\RightR}[1]{}
\providecommand{\TAss}[1]{}
\providecommand{\TRule}[1]{}

\makeatother
\providecommand{\olpath}{.}

% ===================== OLP compatibility layer (2) =====================
% --- second compatibility block: notation and shorthands used by the frozen OLP units ---
\providecommand{\CutCS}{\mathrm{Cut}}
\providecommand{\LeftSubproofLabel}{\relax}
\providecommand{\RightSubproofLabel}{\relax}
\providecommand{\maeh}{\relax}
\providecommand{\cheight}{\relax}
\providecommand{\ande}{\mathrm{and}_E}
\providecommand{\andi}{\mathrm{and}_I}
\providecommand{\ore}{\mathrm{or}_E}
\providecommand{\ori}{\mathrm{or}_I}
\providecommand{\typeof}{\mathord{:}}
\providecommand{\pValue}[1]{#1}
\providecommand{\Sent}{\mathrm{Sent}}
\providecommand{\cardplus}{\mathbin{+}}
\providecommand{\cardtimes}{\mathbin{\cdot}}
\providecommand{\cardexpo}{\mathbin{\hat{\ }}}
\providecommand{\cardnless}{\not<}
\providecommand{\cardle}{\leq}
\providecommand{\funfromto}[3]{#1\colon #2 \to #3}
\providecommand{\funimage}[1]{\operatorname{im}(#1)}
\providecommand{\fdefined}{\mathord{\downarrow}}
\providecommand{\canonord}{\operatorname{ord}}
\providecommand{\ordtimes}{\mathbin{\cdot}}
\providecommand{\ordeq}[2]{#1 = #2}
\providecommand{\ordneq}[2]{#1 \neq #2}
\providecommand{\ordtype}[1]{\operatorname{ord}(#1)}
\providecommand{\ordsucc}{\mathord{+1}}
\providecommand{\ordunion}{\cup}
\providecommand{\onesphere}{S^1}
\providecommand{\rotationsgroup}{\mathrm{SO}(2)}
\providecommand{\equivrep}[1]{[#1]}
\providecommand{\rlexless}{\prec}
\providecommand{\ZFC}{\mathrm{ZFC}}
\providecommand{\ZF}{\mathrm{ZF}}
\providecommand{\Rat}{\mathbb{Q}}
\providecommand{\Real}{\mathbb{R}}
\providecommand{\PosInt}{\mathbb{Z}^{+}}
\providecommand{\Id}{\mathrm{id}}
\providecommand{\Bin}{\mathrm{Bin}}
\providecommand{\pto}{\rightharpoonup}
\providecommand{\dom}{\operatorname{dom}}
\providecommand{\ran}{\operatorname{ran}}
\providecommand{\len}{\operatorname{len}}
\providecommand{\concat}{\mathbin{{}^\frown}}
\providecommand{\card}[1]{\lvert #1 \rvert}
\providecommand{\Pow}{\mathcal{P}}
\providecommand{\Setabs}[2]{\{\, #1 \mid #2 \,\}}
\providecommand{\comp}{\circ}
\providecommand{\Char}{\chi}
\providecommand{\Prf}{\operatorname{Prf}}
\providecommand{\Prov}{\operatorname{Prov}}
\providecommand{\Proj}{\operatorname{proj}}
\providecommand{\Gn}{\operatorname{gn}}
\providecommand{\gn}{\operatorname{gn}}
\providecommand{\Value}[1]{\operatorname{val}(#1)}
\providecommand{\Deduce}{\mathrm{D}}
\providecommand{\Dual}{\mathrm{d}}
\providecommand{\Ax}{\mathrm{Ax}}
\providecommand{\pSat}[3]{#1 \models #2\,[#3]}
\providecommand{\pAssign}{\sigma}
\providecommand{\bforall}{\forall}
\providecommand{\bexists}{\exists}
\providecommand{\liff}{\mathbin{\leftrightarrow}}
\providecommand{\FV}{\operatorname{FV}}
\providecommand{\UnaryInf}{\ensuremath{\mid}}
\providecommand{\BinaryInf}{\ensuremath{\mid}}
\providecommand{\depth}{\operatorname{depth}}
\providecommand{\oliflabeldef}[3]{#1}
\providecommand{\TMblank}{\sqcup}
\providecommand{\TMstroke}{\mid}
\providecommand{\TMstay}{\mathrm{S}}
\providecommand{\TMright}{\mathrm{R}}
\providecommand{\TMleft}{\mathrm{L}}
\providecommand{\TMendtape}{\triangleright}
\providecommand{\umin}{\mu}
\providecommand{\Sum}{\sum}
\providecommand{\citep}[1]{\cite{#1}}
\providecommand{\citet}[1]{\cite{#1}}
\providecommand{\citeauthor}[1]{\cite{#1}}
\providecommand{\citeyear}[1]{\cite{#1}}

% ===================== OLP compatibility layer (3) =====================
% --- third compatibility block: remaining notation seen in the frozen OLP units ---
\providecommand{\iso}{\cong}
\providecommand{\nszero}{\mathrm{0}}
\providecommand{\nssucc}{\mathrm{S}}
\providecommand{\nsplus}{\mathbin{+}}
\providecommand{\nssless}{<}
\providecommand{\nsless}{<}
\providecommand{\PL}{\mathrm{PL}}
\providecommand{\Nec}{\mathord{\Box}}
\providecommand{\RK}{\mathrm{RK}}
\providecommand{\ST}{\mathrm{ST}}
\providecommand{\SSubst}[3]{#1[#2/#3]}
\providecommand{\subst}[3]{#1[#2/#3]}
\providecommand{\tagfalse}{\mathrm{F}}
\providecommand{\cutr}{\mathrm{Cut}}
\providecommand{\cardneq}{\neq}
\providecommand{\issucc}[1]{\operatorname{succ}(#1)}
\providecommand{\tup}[1]{\langle #1 \rangle}
\providecommand{\Var}{\mathrm{Var}}
\providecommand{\Bew}{\operatorname{Bew}}
\providecommand{\Con}{\operatorname{Con}}
\providecommand{\Cons}{\operatorname{Cons}}
\providecommand{\pDeduce}{\mathrm{D}}
\providecommand{\psat}{\models}
\providecommand{\den}[1]{\llbracket #1 \rrbracket}
\providecommand{\Sem}[2]{\llbracket #1 \rrbracket_{#2}}
\providecommand{\DomainOf}[1]{\lvert #1 \rvert}
\providecommand{\Form}{\mathrm{Form}}
\providecommand{\Terms}{\mathrm{Term}}
\providecommand{\Atom}{\mathrm{Atom}}
\providecommand{\FVars}[1]{\operatorname{FV}(#1)}
\providecommand{\fv}[1]{\operatorname{fv}(#1)}
\providecommand{\tuplep}[1]{\langle #1 \rangle}
\providecommand{\restrict}[2]{#1 \upharpoonright #2}
\providecommand{\sto}{\rightarrow}
\providecommand{\FmlaSet}[1]{\mathrm{Fmla}_{#1}}
\providecommand{\Thm}{\mathrm{Thm}}
\providecommand{\Eval}{\operatorname{val}}
\providecommand{\blank}{\sqcup}
\providecommand{\blankTape}{\sqcup}
\providecommand{\Consis}{\operatorname{Con}}
\providecommand{\drs}{\mathrm{d}}
\providecommand{\VRule}{\mid}
\providecommand{\Bet}{\mathrm{B}}
\providecommand{\Seg}[1]{\mathrm{Seg}_{#1}}
\providecommand{\Enums}{\mathrm{En}}
\providecommand{\Z}{\mathbb{Z}}
\providecommand{\N}{\mathbb{N}}
\providecommand{\abs}[1]{\lvert #1 \rvert}
\providecommand{\floor}[1]{\lfloor #1 \rfloor}
\providecommand{\ceil}[1]{\lceil #1 \rceil}

% ===================== OLP compatibility layer (4) =====================
% --- fourth compatibility block: harvested from the per-candidate compile logs ---
\providecommand{\beredpar}{\mathrel{\rightarrow_\beta}}
\makeatletter
\def\bered{\@ifnextchar[{\bered@i}{\mathrel{\rightarrow_\beta}}}
\def\bered@i[#1]{\mathrel{\rightarrow_\beta}}
\def\bred{\@ifnextchar[{\bred@i}{\mathrel{\rightarrow_\beta}}}
\def\bred@i[#1]{\mathrel{\rightarrow_\beta}}
\makeatother
\providecommand{\beredone}{\mathrel{=_\beta}}
\providecommand{\bredone}{\mathrel{=_\beta}}
\providecommand{\xredone}{\mathrel{=_\beta}}
\providecommand{\aconvone}{\mathrel{=_\alpha}}
\providecommand{\cardless}{<}
\providecommand{\tagrefs}[1]{#1}
\providecommand{\ordplus}{\mathbin{+}}
\providecommand{\funrestrictionto}[2]{#1 \upharpoonright #2}
\providecommand{\fundefined}{\mathord{\uparrow}}
\providecommand{\Trm}{\mathrm{Trm}}
\providecommand{\OProv}{\operatorname{Prov}}
\providecommand{\OPrf}{\operatorname{Prf}}
\providecommand{\aeq}{\mathrel{=}}
\providecommand{\mClass}[1]{[#1]}
\providecommand{\equivclass}[1]{[#1]}
\providecommand{\pheight}{\operatorname{ht}}
\providecommand{\phy}{\operatorname{ht}}
\providecommand{\Cut}{\mathrm{Cut}}
\providecommand{\QuantRank}{\operatorname{qr}}
\providecommand{\elemequiv}{\mathrel{\equiv}}
\providecommand{\PIso}{\mathrm{PIso}}
\providecommand{\Theory}{\mathrm{Th}}
\providecommand{\supstrict}{\subsetneq}
\providecommand{\trcl}{\operatorname{trcl}}
\providecommand{\cref}[1]{\ref{#1}}
\providecommand{\Discharge}{\mathrm{Dis}}
\providecommand{\Weakening}{\mathrm{Wk}}
\providecommand{\ZFminus}{\mathrm{ZF}^{-}}
\makeatletter
\renewcommand{\Dual}{\@ifnextchar[{\Dual@i}{\ifmmode\mathrm{d}\else$\mathrm{d}$\fi}}
\def\Dual@i[#1]{\ifmmode\mathrm{d}[#1]\else$\mathrm{d}[#1]$\fi}
\makeatother
\providecommand{\RightR}{\mathrm{R}}
\providecommand{\LeftR}{\mathrm{L}}
\providecommand{\Intro}{\mathrm{I}}
\providecommand{\Axiom}{\mathrm{Ax}}
\providecommand{\formula}[1]{#1}
\providecommand{\stared}{\ast}
\providecommand{\oset}[2]{\{\, #1 \mid #2 \,\}}
\providecommand{\setof}[2]{\{\, #1 \mid #2 \,\}}
\providecommand{\pow}[1]{\mathcal{P}(#1)}
\providecommand{\ordpair}[2]{\langle #1, #2 \rangle}
\providecommand{\ordtriple}[3]{\langle #1, #2, #3 \rangle}
\providecommand{\union}{\cup}
\providecommand{\inter}{\cap}
\providecommand{\seq}[1]{\langle #1 \rangle}
\providecommand{\SuchThat}{\mid}
\providecommand{\where}{\mid}
\providecommand{\TVar}{\mathrm{TVar}}
\providecommand{\Bool}{\mathrm{Bool}}
\providecommand{\Val}{\operatorname{Val}}
\providecommand{\TL}[1]{\mathrm{TL}}
\providecommand{\xstared}{\ast}

% ===================== OLP compatibility layer (5) =====================
% --- fifth compatibility block: rescues the quarantined OLP candidates ---
\providecommand{\ident}{\operatorname{id}}
\providecommand{\setrank}{\operatorname{rank}}
\providecommand{\cardfont}{\mathfrak{c}}
\providecommand{\Complement}{\overline}
\providecommand{\QR}{\mathrm{QR}}
\providecommand{\becd}{\ensuremath{\ulcorner}}
\providecommand{\disjointsum}{\mathbin{\oplus}}
\providecommand{\cardeq}[2]{#1 = #2}
\providecommand{\xred}{\mathrel{\rightarrow}}
% --- sixth compatibility block: names reported by real compile logs ---
\providecommand{\ORefut}{\operatorname{Refut}}
\providecommand{\ORProv}{\operatorname{Prov}}
\providecommand{\cutrank}{\operatorname{cr}}
\providecommand{\Refut}{\operatorname{Refut}}
\providecommand{\DischargeRule}{\mathrm{Dis}}
\providecommand{\Contraction}{\mathrm{C}}
\providecommand{\bottomAlignProof}{\relax}
\providecommand{\FalseCl}{\bot}
\providecommand{\red}{\mathrel{\rightarrow}}
\providecommand{\varAssign}{\sigma}
\providecommand{\defiff}{\mathrel{\leftrightarrow}}
\providecommand{\FalseInt}{\bot}
\providecommand{\fn}{\operatorname{fn}}
\providecommand{\dotted}{\cdot}

% --- tag-macro semantic fix (2026-10-13, final): \iftag CHOOSES a branch and the OLP proof body is
% always the SECOND group (\iftag{tag}{body}, {tag}{body}{c}, {tag}{body}{c}{d}); a definition that
% prints the tag literal drops the proof body. Bodies contain blank lines, hence \long.
\makeatletter
\providecommand{\iftag}{}
\long\def\iftag#1#2{\@ifnextchar\bgroup{\iftag@three{#2}}{\iftag@stop{#2}}}
\long\def\iftag@three#1#2{\iftag@stop{#1}\@ifnextchar\bgroup{\iftag@eat}{\relax}}
\long\def\iftag@eat#1{}
\long\def\iftag@stop#1{#1}
\long\def\oliflabeldef#1#2{\@ifnextchar\bgroup{\@gobble}{\iftag@stop{#2}}}
\providecommand{\tagrefs}[1]{}
\providecommand{\settag}[1]{}
\makeatother
\begin{document}

% ==== OLP file: content/model-theory/models-of-arithmetic/computable-models.tex ====
% Part: model-theory
% Chapter: models-of-arithmetic
% Section: computable-models



\olfileid{mod}{mar}{cmp}
\section{Computable Models of Arithmetic}

\begin{explain}
The standard model~$\Struct{N}$ has two nice features. Its domain is
the natural numbers~$\Nat$, i.e., its elements are just the kinds of
things we want to talk about using the language of arithmetic, and the
standard numeral~$\num{n}$ actually picks out~$n$. The other nice
feature is that the interpretations of the non-logical symbols
of~$\Lang{L_A}$ are all \emph{computable}. The successor, addition,
and multiplication functions which serve as $\Assign{\prime}{N}$,
$\Assign{+}{N}$, and $\Assign{\times}{N}$ are computable functions of
numbers. (Computable by Turing machines, or definable by primitive
recursion, say.) And the less-than relation on~$\Struct{N}$, i.e.,
$\Assign{<}{N}$, is decidable.

Non-standard models of arithmetical theories such as $\Th{Q}$ and
$\Th{PA}$ must contain non-standard elements. Thus their domains
typically include \olup{element}s in addition to~$\Nat$. However, any
countable \olup{structure} can be built on any \olup{denumerable} set,
including~$\Nat$. So there are also non-standard models with
domain~$\Nat$. In such models~$\Struct{M}$, of course, at least some
numbers cannot play the roles they usually play, since some~$k$ must
be different from~$\Value{\num{n}}{M}$ for all~$n \in \Nat$.
\end{explain}

\begin{defn}
!!^a{structure}~$\Struct{M}$ for $\Lang{L_A}$ is \emph{computable} iff
  $\Domain{M} = \Nat$ and $\Assign{\prime}{M}$, $\Assign{+}{M}$,
  $\Assign{\times}{M}$ are computable functions and $\Assign{<}{M}$ is
  a decidable relation.
\end{defn}

\begin{ex}\ollabel{ex:comp-model-q}
Recall the structure $\Struct{K}$ from \olref[mdq]{ex:model-K-of-Q}.
Its domain was $\Domain{K} = \Nat
\cup \{a\}$ and interpretations
\begin{align*}
  \Assign{\Obj{0}}{K} & = 0\\
  \Assign{\prime}{K}(x) & =
  \begin{cases}
    x+1 & \text{if $x\in \Nat$}\\
    a & \text{if $x = a$}
  \end{cases}\\
  \Assign{+}{K}(x, y) & =
  \begin{cases}
    x+y & \text{if $x$, $y \in\Nat$}\\
    a & \text{otherwise}
  \end{cases}\\
  \Assign{\times}{K}(x, y) & =
  \begin{cases}
    xy & \text{if $x$, $y \in\Nat$}\\
    0 & \text{if $x=0$ or $y=0$}\\
    a & \text{otherwise}\\
  \end{cases}\\
  \Assign{<}{K} & =
  \Setabs{\tuple{x,y}}{x, y \in \Nat \text{ and } x<y} \cup
  \Setabs{\tuple{x,a}}{n \in \Domain{K}}
\end{align*}
But $\Domain{K}$ is \olup{denumerable} and so is equinumerous
with~$\Nat$. For instance, $g\colon \Nat \to \Domain{K}$ with $g(0) =
a$ and $g(n) = n+1$ for $n>0$ is \olup{bijection}.  We can turn it into
an isomorphism between a new model~$\Struct{K'}$ of~$\Th{Q}$ and
$\Struct{K}$.  In $\Struct{K'}$, we have to assign different functions
and relations to the symbols of~$\Lang{L_A}$, since different
\olup{element}s of~$\Nat$ play the roles of standard and non-standard
numbers.

Specifically, $0$ now plays the role of~$a$, not of the smallest
standard number. The smallest standard number is now~$1$. So we assign
$\Assign{\Obj{0}}{K'} = 1$. The successor function is also different
now: given a standard number, i.e., an $n > 0$, it still returns
$n+1$. But $0$ now plays the role of~$a$, which is its own
successor. So $\Assign{\prime}{K'}(0) = 0$.  For addition and
multiplication we likewise have
\begin{align*}
\Assign{+}{K'}(x, y) & =
  \begin{cases}
    x+y-1 & \text{if $x$, $y >0$}\\
    0 & \text{otherwise}
  \end{cases}\\
  \Assign{\times}{K'}(x, y) & =
  \begin{cases}
    1 & \text{if $x = 1$ or $y = 1$}\\
    xy - x - y + 2 & \text{if $x$, $y > 1$}\\
    0 & \text{otherwise}\\
  \end{cases}
\end{align*}
And we have $\tuple{x, y} \in \Assign{<}{K'}$ iff $x < y$ and $x > 0$
and $y > 0$, or if $y = 0$.

All of these functions are computable functions of natural numbers and
$\Assign{<}{K'}$ is a decidable relation on~$\Nat$---but they are not
the same functions as successor, addition, and multiplication
on~$\Nat$, and $\Assign{<}{K'}$ is not the same relation as~$<$
on~$\Nat$.
\end{ex}

\begin{prob}
Give \olup{structure}~$\Struct{L'}$ with $\Domain{L'} = \Nat$ isomorphic
to~$\Struct{L}$ of \olref[mod][mar][mdq]{ex:model-L-of-Q}.
\end{prob}

\begin{explain}
\Olref{ex:comp-model-q} shows that $\Th{Q}$ has computable
non-standard models with domain~$\Nat$.  However, the following result
shows that this is not true for models of~$\Th{PA}$ (and thus also for
models of~$\Th{TA}$).
\end{explain}

\begin{thm}[Tennenbaum's Theorem]
$\Struct{N}$ is the only computable model of~$\Th{PA}$.
\end{thm}
  

% ==== OLP file: content/model-theory/models-of-arithmetic/introduction.tex ====
% Part: model-theory
% Chapter: models-of-arithmetic
% Section: introduction



\olfileid{mod}{mar}{int}
\section{Introduction}

The \emph{standard model} of arithmetic is the
\olup{structure}~$\Struct{N}$ with $\Domain{N} = \Nat$ in which
$\Obj{0}$, $\prime$, $+$, $\times$, and $<$ are interpreted as you
would expect. That is, $\Obj{0}$ is $0$, $\prime$ is the successor
function, $+$ is interpreted as addition and $\times$ as multiplication
of the numbers in~$\Nat$. Specifically,
\begin{align*}
  \Assign{\Obj{0}}{N} & = 0\\
  \Assign{\prime}{N}(n) & = n + 1\\
  \Assign{+}{N}(n, m) & = n + m\\
  \Assign{\times}{N}(n, m) & = nm
\end{align*}
Of course, there are structures for $\Lang{L_A}$ that have domains
other than~$\Nat$. For instance, we can take $\Struct{M}$ with domain
$\Domain{M} = \{a\}^*$ (the finite sequences of the single
symbol~$a$, i.e., $\emptyset$, $a$, $aa$, $aaa$, \dots), and
interpretations
\begin{align*}
  \Assign{\Obj{0}}{M} & = \emptyset\\
  \Assign{\prime}{M}(s) & = s \concat a\\
  \Assign{+}{M}(n, m) & = a^{n + m}\\
  \Assign{\times}{M}(n, m) & = a^{nm}
\end{align*}
These two structures are ``essentially the same'' in the sense that
the only difference is the \olup{element}s of the \olup{domain}s but not how
the \olup{element}s of the \olup{domain}s are related among each other by
the interpretation functions. We say that the two \olup{structure}s are
\emph{isomorphic}.

It is an easy consequence of the compactness theorem that any theory
true in~$\Struct{N}$ also has models that are not isomorphic
to~$\Struct{N}$.  Such structures are called \emph{non-standard}.  The
interesting thing about them is that while the \olup{element}s of a
standard model (i.e., $\Struct{N}$, but also all \olup{structure}s
isomorphic to it) are exhausted by the values of the standard
numerals~$\num{n}$, i.e.,
\[
\Domain{N} = \Setabs{\Value{\num{n}}{N}}{n \in \Nat}
\]
that isn't the case in non-standard models: if $\Struct{M}$ is
non-standard, then there is at least one $x \in \Domain{M}$ such that
$x \neq \Value{\num{n}}{M}$ for all~$n$.

These non-standard elements are pretty neat: they are ``infinite
natural numbers.'' But their existence also explains, in a sense, the
incompleteness phenomena.  Consider an example, e.g., the consistency
statement for Peano arithmetic, $\OCon[\Th{PA}]$, i.e., $\lnot
\lexists[x][\OPrf[\Th{PA}](x, \gn{\lfalse})]$. Since $\Th{PA}$ neither
proves $\OCon[\Th{PA}]$ nor $\lnot \OCon[\Th{PA}]$, either can be
consistently added to $\Th{PA}$. Since $\Th{PA}$ is consistent,
$\Sat{N}{\OCon[\Th{PA}]}$, and consequently $\Sat/{N}{\lnot
  \OCon[\Th{PA}]}$.  So $\Struct{N}$ is \emph{not} a model of $\Th{PA}
\cup \{\lnot \OCon[\Th{PA}]\}$, and all its models must be
nonstandard. Models of $\Th{PA} \cup \{\lnot \OCon[\Th{PA}]\}$ must
contain some \olup{element} that serves as the witness that makes
$\lexists[x][\OPrf[\Th{PA}](\gn{\lfalse})]$ true, i.e., a G\"odel
number of \olup{derivation} of a contradiction from~$\Th{PA}$.  Such
\olup{element} can't be standard---since $\Th{PA} \Proves \lnot
\OPrf[\Th{PA}](\num{n}, \gn{\lfalse})$ for every~$n$.


% ==== OLP file: content/model-theory/models-of-arithmetic/models-of-arithmetic.tex ====
% Part: model-theory
% Chapter: models-of-arithmetic



\olchapter{mod}{mar}{Models of Arithmetic}

\olimport{introduction}

\olimport{standard-models}

\olimport{non-standard-models}

\olimport{models-of-q}

\olimport{models-of-pa}

\olimport{computable-models}

\OLEndChapterHook


% ==== OLP file: content/model-theory/models-of-arithmetic/models-of-pa.tex ====
% Part: first-order-logic
% Chapter: models-of-arithmetic
% Section: models-of-pa



\olfileid{mod}{mar}{mpa}
\section{Models of $\Th{PA}$}

\begin{explain}
Any non-standard model of~$\Th{TA}$ is also one of~$\Th{PA}$. We know
that non-standard models of~$\Th{TA}$ and hence of~$\Th{PA}$ exist. We
also know that such non-standard models contain non-standard
``numbers,'' i.e., \olup{element}s of the domain that are ``beyond'' all
the standard ``numbers.''  But how are they arranged? How many are
there? We've seen that models of the weaker theory~$\Th{Q}$ can
contain as few as a single non-standard number. But these simple
\olup{structure}s are not models of $\Th{PA}$ or $\Th{TA}$.

The key to understanding the structure of models of $\Th{PA}$ or
$\Th{TA}$ is to see what facts are \olup{derivable} in these theories.
For instance, already $\Th{PA}$ proves that $\lforall[x][\eq/[x][x']]$
and $\lforall[x][\lforall[y][\eq[(x+y)][(y+x)]]]$, so this rules out
simple structures (in which these \olup{sentence}s are false) as models
of~$\Th{PA}$.

Suppose~$\Struct{M}$ is a model of~$\Th{PA}$. Then if $\Th{PA} \Proves
A$, $\Sat{M}{A}$.  Let's again use $\nszero$ for
$\Assign{\Obj{0}}{M}$, $\nssucc$ for $\Assign{\prime}{M}$, $\nsplus$
for $\Assign{+}{M}$, $\nstimes$ for $\Assign{\times}{M}$, and
$\nsless$ for $\Assign{<}{M}$. Any \olup{sentence}~$A$ then states some
condition about $\nszero$, $\nssucc$, $\nsplus$, $\nstimes$, and
$\nsless$, and if $\Sat{M}{A}$ that condition must be satisfied. For
instance, if $\Sat{M}{Q_1}$, i.e.,
$\Sat{M}{\lforall[x][\lforall[y][(\eq[x'][y'] \lif \eq[x][y])]]}$,
then $\nssucc$ must be \olup{injective}.
\end{explain}

\begin{prop}
In $\Struct{M}$, $\nsless$ is a linear strict order, i.e., it
satisfies:
\begin{enumerate}
\item Not $x \nsless x$ for any~$x \in \Domain{M}$.
\item If $x \nsless y$ and $y \nsless z$ then $x \nsless z$.
\item For any $x \neq y$, $x \nsless y$ or $y \nsless x$
\end{enumerate}
\end{prop}

\begin{proof}
$\Th{PA}$ proves:
\begin{enumerate}
\item $\lforall[x][\lnot x < x]$
\item $\lforall[x][\lforall[y][\lforall[z][((x < y \land y < z) \lif x < z)]]]$
\item $\lforall[x][\lforall[y][((x < y \lor y < x) \lor \eq[x][y]))]]$
\end{enumerate}
\end{proof}

\begin{prop}
\ollabel{prop:M-discrete} $\nszero$ is the least \olup{element}
of~$\Domain{M}$ in the $\nsless$-ordering. For any $x$, $x \nsless
x^\nssucc$, and $x^\nssucc$ is the $\nsless$-least \olup{element} with
that property.  For any $x$, there is a unique $y$ such that
$y^\nssucc = x$. (We call $y$ the ``predecessor'' of~$x$
in~$\Struct{M}$, and denote it by~$^\nssucc x$.)
\end{prop}

\begin{proof}
Exercise.  
\end{proof}

\begin{prob}
Find \olup{sentence}s in~$\Lang{L_A}$ \olup{derivable} in~$\Th{PA}$ (and
hence true in~$\Struct{N}$) which guarantee the properties of
$\nszero$, $\nssucc$, and $\nsless$ in
\olref[mod][mar][mpa]{prop:M-discrete}
\end{prob}

\begin{prop}
All standard \olup{element}s of~$\Struct{M}$ are less than (according
to~$\nsless$) all non-standard \olup{element}s.
\end{prop}

\begin{proof}
We'll use $n$ as short for $\Value{\num{n}}{M}$, a standard
\olup{element} of~$\Struct{M}$.  Already $\Th{Q}$ proves that, for any~$n
\in \Nat$, $\lforall[x][(x < \num{n}' \lif (\eq[x][\num{0}] \lor
  \eq[x][\num{1}] \lor \dots \lor \eq[x][\num{n}]))]$. There are no
\olup{element}s that are $\nsless \nszero$. So if $n$ is standard and $x$
is non-standard, we cannot have $x \nsless n$. By definition, a
non-standard element is one that isn't $\Value{\num{n}}{M}$ for any~$n
\in \Nat$, so $x \neq n$ as well. Since $\nsless$ is a linear order,
we must have $n \nsless x$.
\end{proof}

\begin{prop}
Every nonstandard \olup{element}~$x$ of~$\Domain{M}$ is an element of the subset
\[
\dots ^{\nssucc\nssucc\nssucc}x \nsless ^{\nssucc\nssucc}x \nsless
^{\nssucc}x \nsless x \nsless x^{\nssucc} \nsless x^{\nssucc\nssucc}
\nsless x^{\nssucc\nssucc\nssucc} \nsless \dots
\]
We call this subset the \emph{block of~$x$} and write it as $[x]$. It
has no least and no greatest \olup{element}. It can be characterized as
the set of those $y \in \Domain{M}$ such that, for some standard~$n$,
$x \nsplus n = y$ or $y \nsplus n = x$.
\end{prop}

\begin{proof}
Clearly, such a set~$[x]$ always exists since every \olup{element}~$y$ of
$\Domain{M}$ has a unique successor~$y^\nssucc$ and unique
predecessor~$^\nssucc y$. For successive \olup{element}s $y$, $y^\nssucc$
we have $y \nsless y^\nssucc$ and $y^\nssucc$ is the $\nsless$-least
\olup{element} of~$\Domain{M}$ such that $y$ is $\nsless$-less than it.
Since always $^\nssucc y \nsless y$ and $y \nsless y^\nssucc$, $[x]$
has no least or greatest \olup{element}. If $y \in [x]$ then $x \in [y]$,
for then either $y^{\nssucc\dots\nssucc} = x$ or
$x^{\nssucc\dots\nssucc} = y$.  If $y^{\nssucc\dots\nssucc} = x$ (with
$n$ $\nssucc$'s), then $y \nsplus n = x$ and conversely, since
$\Th{PA} \Proves \lforall[x][\eq[x^{\prime\dots\prime}][(x +
    \num{n})]]$ (if $n$ is the number of~$\prime$'s).
\end{proof}

\begin{prop}
If $[x] \neq [y]$ and $x \nsless y$, then for any $u \in [x]$ and any
$v \in [y]$, $u \nsless v$.
\end{prop}

\begin{proof}
Note that $\Th{PA} \Proves \lforall[x][\lforall[y][(x < y \lif (x' < y
    \lor x' = y))]]$. Thus, if $u \nsless v$, we also have $u \nsplus
n^\nssucc \nsless v$ for any~$n$ if $[u] \neq [v]$.

Any $u \in [x]$ is $\nsless y$: $x \nsless y$ by assumption. If $u
\nsless x$, $u \nsless y$ by transitivity. And if $x \nsless u$ but $u
\in [x]$, we have $u = x\nsplus n^\nssucc$ for some~$n$, and so $u
\nsless y$ by the fact just proved.
  
Now suppose that $v \in [y]$ is $\nsless y$, i.e., $v \nsplus
m^\nssucc = y$ for some standard~$m$. This rules out $v \nsless x$,
otherwise $y = v \nsplus m^\nssucc \nsless x$. Clearly also, $x \neq
v$, otherwise $x \nsplus m^\nssucc = v \nsplus m^\nssucc = y$ and we
would have $[x] = [y]$. So, $x \nsless v$. But then also $x \nsplus
n^\nssucc \nsless v$ for any~$n$. Hence, if $x \nsless u$ and $u \in
[x]$, we have $u \nsless v$. If $u \nsless x$ then $u \nsless v$ by
transitivity.

Lastly, if $y \nsless v$, $u \nsless v$ since, as we've shown, $u
\nsless y$ and $y \nsless v$.
\end{proof}

\begin{cor}
If $[x] \neq [y]$, $[x] \cap [y] = \emptyset$.
\end{cor}

\begin{proof}
Suppose $z \in [x]$ and $x \nsless y$. Then $z \nsless u$ for all $u
\in [y]$. If $z \in [y]$, we would have $z \nsless z$. Similarly if $y
\nsless x$.
\end{proof}

\begin{explain}
This means that the blocks themselves can be ordered in a way that
respects $\nsless$: $[x] \nsless [y]$ iff $x \nsless y$, or,
equivalently, if $u \nsless v$ for any $u \in [x]$ and $v \in [y]$.
Clearly, the standard block $[0]$ is the least block. It intersects
with no non-standard block, and no two non-standard blocks intersect
either. Specifically, you cannot ``reach'' a different block by taking
repeated successors or predecessors.
\end{explain}

\begin{prop}
If $x$ and $y$ are non-standard, then $x \nsless x \nsplus y$ and $x
\nsplus y \notin [x]$.
\end{prop}

\begin{proof}
If $y$ is nonstandard, then $y \neq \nszero$. $\Th{PA} \Proves
\lforall[x][(y \neq \Obj{0} \lif x < (x+y))]$.  Now suppose $x
\nsplus y \in [x]$. Since $x \nsless x \nsplus y$, we would have $x
\nsplus n^\nssucc = x \nsplus y$. But $\Th{PA} \Proves
\lforall[x][\lforall[y][\lforall[z][(\eq[(x+y)][(x+z)] \lif y = z)]]]$
(the cancellation law for addition). This would mean $y = n^\nssucc$
for some standard~$n$; but $y$ is assumed to be non-standard.
\end{proof}

\begin{prop}
There is no least non-standard block.
\end{prop}

\begin{proof}
$\Th{PA} \Proves \lforall[x][\lexists[y][(\eq[(y+y)][x] \lor
      \eq[(y+y)'][x])]]$, i.e., that every $x$ is divisible by~$2$
  (possibly with remainder~$1$). If $x$ is non-standard, so is~$y$. By
  the preceding proposition, $y \nsless y \nsplus y$ and $y \nsplus y
  \notin [y]$. Then also $y \nsless (y \nsplus y)^\nssucc$ and $(y
  \nsplus y)^\nssucc \notin [y]$. But $x = y \nsplus y$ or $x = (y
  \nsplus y)^\nssucc$, so $y \nsless x$ and $y \notin [x]$.
\end{proof}

\begin{prop}
There is no largest block.
\end{prop}

\begin{proof}
Exercise.
\end{proof}

\begin{prob}
Show that in a non-standard model of~$\Th{PA}$, there is no largest
block.
\end{prob}

\begin{prop}
\ollabel{prop:blocks-dense}
The ordering of the blocks is dense. That is, if $x \nsless y$ and
$[x] \neq [y]$, then there is a block $[z]$ distinct from both that is
between them.
\end{prop}

\begin{proof}
Suppose $x \nsless y$. As before, $x \nsplus y$ is divisible by two
(possibly with remainder): there is a $z \in \Domain{M}$ such that
either $x \oplus y = z \oplus z$ or $x \oplus y = (z \oplus
z)^\nssucc$. The element $z$ is the ``average'' of $x$ and~$y$, and $x
\nsless z$ and $z \nsless y$.
\end{proof}

\begin{prob}
Write out a detailed proof of
\olref[mod][mar][mpa]{prop:blocks-dense}. Which \olup{sentence} must
$\Th{PA}$ \olup{derive} in order to guarantee the existence of~$z$? Why
is $x \nsless z$ and $z \nsless y$, and why is $[x] \neq [z]$ and $[z]
\neq [y]$?
\end{prob}

\begin{explain}
The non-standard blocks are therefore ordered like the rationals: they
form \olup{denumerable} dense linear ordering without endpoints.  One can show
that any two such \olup{denumerable} orderings are isomorphic. It follows
that for any two \olup{enumerable} non-standard models $\Struct{M}_1$ and
$\Struct{M_2}$ of true arithmetic, their reducts to the language
containing $<$ and $=$ only are isomorphic. Indeed, an isomorphism $h$
can be defined as follows: the standard parts of $\Struct{M_1}$ and
$\Struct{M_2}$ are isomorphic to the standard model $\Struct{N}$ and
hence to each other. The blocks making up the non-standard part are
themselves ordered like the rationals and therefore isomorphic; an
isomorphism of the blocks can be extended to an isomorphism
\emph{within} the blocks by matching up arbitrary elements in each,
and then taking the image of the successor of $x$ in $\Struct{M_1}$ to
be the successor of the image of $x$ in $\Struct{M_2}$. Note that it
does \emph{not} follow that $\mathfrak{M}_1$ and $\mathfrak{M}_2$ are
isomorphic in the full language of arithmetic (indeed, isomorphism is
always relative to \olup{language}), as there are non-isomorphic ways to
define addition and multiplication over $\Domain{M_1}$ and
$\Domain{M_2}$. (This also follows from a famous theorem due to Vaught
that the number of countable models of a complete theory cannot be~2.)
\end{explain}

% ==== OLP file: content/model-theory/models-of-arithmetic/models-of-q.tex ====
% Part: model-theory
% Chapter: models-of-arithmetic
% Section: non-standard-models



\olfileid{mod}{mar}{mdq}
\section{Models of $\Th{Q}$}

\begin{explain}
We know that there are non-standard \olup{structure}s that make the same
\olup{sentence}s true as~$\Struct{N}$ does, i.e., is a model
of~$\Th{TA}$. Since $\Sat{N}{\Th{Q}}$, any model of~$\Th{TA}$ is also
a model of~$\Th{Q}$. $\Th{Q}$ is much weaker than~$\Th{TA}$, e.g.,
$\Th{Q} \Proves/ \lforall[x][\lforall[y][\eq[(x +
      y)][(y+x)]]]$. Weaker theories are easier to satisfy: they have
more models. E.g., $\Th{Q}$ has models which make
$\lforall[x][\lforall[y][\eq[(x + y)][(y+x)]]]$ false, but those
cannot also be models of~$\Th{TA}$, or $\Th{PA}$ for that matter.
Models of $\Th{Q}$ are also relatively simple: we can specify them
explicitly.
\end{explain}

\begin{ex}
\ollabel{ex:model-K-of-Q}
Consider the \olup{structure}~$\Struct{K}$ with domain $\Domain{K} = \Nat
\cup \{a\}$ and interpretations
\begin{align*}
  \Assign{\Obj{0}}{K} & = 0\\
  \Assign{\prime}{K}(x) & =
  \begin{cases}
    x+1 & \text{if $x\in \Nat$}\\
    a & \text{if $x = a$}
  \end{cases}\\
  \Assign{+}{K}(x, y) & =
  \begin{cases}
    x+y & \text{if $x$, $y \in\Nat$}\\
    a & \text{otherwise}
  \end{cases}\\
  \Assign{\times}{K}(x, y) & =
  \begin{cases}
    xy & \text{if $x$, $y \in\Nat$}\\
    0 & \text{if $x = 0$ or $y = 0$}\\
    a & \text{otherwise}\\
  \end{cases}\\
  \Assign{<}{K} & =
  \Setabs{\tuple{x,y}}{x, y \in \Nat \text{ and } x<y} \cup
  \Setabs{\tuple{x,a}}{x \in \Domain{K}}
\end{align*}
To show that $\Sat{K}{\Th{Q}}$ we have to verify that all axioms
of~$\Th{Q}$ are true in~$\Struct{K}$.  For convenience, let's write
$x^\nssucc$ for $\Assign{\prime}{K}(x)$ (the ``successor'' of $x$
in~$\Struct{K}$), $x \nsplus y$ for $\Assign{+}{K}(x, y)$ (the ``sum''
of $x$ and $y$ in~$\Struct{K}$, $x \nstimes y$ for
$\Assign{\times}{K}(x, y)$ (the ``product'' of $x$ and~$y$
in~$\Struct{K}$), and $x \nsless y$ for $\tuple{x,y} \in
\Assign{<}{K}$. With these abbreviations, we can give the operations
in~$\Struct{K}$ more perspicuously as
\[
\begin{array}{c|c}
  x & x^\nssucc \\
  \hline
  n & n+1 \\
  a & a
\end{array}
\qquad
\begin{array}{c|ccc}
  x \nsplus y & 0 & m & a \\
  \hline
  0 & 0 & m & a \\
  n & n & n+m & a \\
  a & a & a & a \\
\end{array}
\qquad
\begin{array}{c|ccc}
  x \nstimes y & 0 & m & a \\
  \hline
  0 & 0 & 0 & 0 \\
  n & 0 & nm & a \\
  a & 0 & a & a \\
\end{array}
\]
We have $n \nsless m$ iff $n<m$ for $n$, $m \in \Nat$ and $x \nsless
a$ for all~$x \in \Domain{K}$.

$\Sat{K}{\lforall[x][\lforall[y][(\eq[x'][y'] \lif \eq[x][y])]]}$
since $\nssucc$ is \olup{injective}. $\Sat{K}{\lforall[x][\eq/[\Obj
0][x']]}$ since $0$~is not a $\nssucc$-successor in~$\Struct{K}$.
$\Sat{K}{\lforall[x][(\eq[x][\Obj 0] \lor \lexists[y][\eq[x][y']])]}$
since for every $n>0$, $n = (n-1)^\nssucc$, and $a = a^\nssucc$.

$\Sat{K}{\lforall[x][\eq[(x + \Obj 0)][x]]}$ since $n \nsplus 0 = n+0
= n$, and $a\nsplus 0 = a$ by definition of~$\nsplus$.
$\Sat{K}{\lforall[x][\lforall[y][\eq[(x + y')][(x + y)']]]}$ is a bit
trickier.  If $n$, $m$ are both standard, we have:
\begin{align*}
(n \nsplus m^\nssucc) & = (n+(m+1)) = (n+m)+1 = (n \nsplus m)^\nssucc  
\intertext{since $\nsplus$ and $^\nssucc$ agree with $+$ and $\prime$ on
  standard numbers.  Now suppose $x \in \Domain{K}$. Then}
(x \nsplus a^\nssucc) & = (x \nsplus a) = a = a^\nssucc = (x \nsplus a)^\nssucc
\intertext{The remaining case is if $y \in \Domain{K}$ but $x =
  a$. Here we also have to distinguish cases according to whether $y =
  n$ is standard or $y = b$:}
(a \nsplus n^\nssucc) & = (a \nsplus (n+1)) = a = a^\nssucc = (a \nsplus n)^\nssucc\\
(a \nsplus a^\nssucc) & = (a \nsplus a) = a = a^\nssucc = (a \nsplus a)^\nssucc
\end{align*}
This is of course a bit more detailed than needed. For instance, since
$a \nsplus z = a$ whatever $z$ is, we can immediately conclude $a \nsplus
a^\nssucc = a$. The remaining axioms can be verified the same way.

$\Struct{K}$ is thus a model of~$\Th{Q}$. Its ``addition''~$\nsplus$
is also commutative. But there are other sentences true
in~$\Struct{N}$ but false in~$\Struct{K}$, and vice versa. For
instance, $a \nsless a$, so $\Sat{K}{\lexists[x][x < x]}$ and
$\Sat/{K}{\lforall[x][\lnot x<x]}$. This shows that $\Th{Q} \Proves/
\lforall[x][\lnot x < x]$.
\end{ex}

\begin{prob}
Prove that $\Struct{K}$ from \olref[mod][mar][mdq]{ex:model-K-of-Q}
satisfies the remaining axioms of~$\Th{Q}$,
\begin{align*}
  & \lforall[x][\eq[(x \times \Obj 0)][\Obj 0]] \tag{$Q_6$}\\
  & \lforall[x][\lforall[y][\eq[(x \times y')][((x \times y) + x)]]] \tag{$Q_7$}\\
  & \lforall[x][\lforall[y][(x < y \liff \lexists[z][\eq[(z' + x)][y]])]] \tag{$Q_8$}
\end{align*}
Find \olup{sentence} only involving~$\prime$ true in~$\Struct{N}$ but
false in~$\Struct{K}$.
\end{prob}

\begin{ex}
\ollabel{ex:model-L-of-Q} Consider the \olup{structure}~$\Struct{L}$ with
domain $\Domain{L} = \Nat \cup \{a, b\}$ and interpretations
$\Assign{\prime}{L} = \nssucc$, $\Assign{+}{L} = \nsplus$ given by
\[
\begin{array}{c|c}
  x & x^\nssucc \\
  \hline
  n & n+1 \\
  a & a\\
  b & b
\end{array}
\qquad
\begin{array}{c|ccc}
  x \nsplus y & m & a & b\\
  \hline
  n & n+m & b & a\\
  a & a & b & a\\
  b & b & b & a
\end{array}
\]
Since $\nssucc$ is \olup{injective}, $0$ is not in its range, and every
$x \in \Domain{L}$ other than~$0$ is, axioms $Q_1$--$Q_3$ are true
in~$\Struct{L}$. For any $x$, $x \nsplus 0 = x$, so $Q_4$ is true as
well.  For $Q_5$, consider $x \nsplus y^\nssucc$ and $(x \nsplus
y)^\nssucc$. They are equal if $x$ and $y$ are both standard, since
then $\nssucc$ and $\nsplus$ agree with $\prime$ and $+$. If $x$ is
non-standard, and $y$ is standard, we have $x \nsplus y^\nssucc = x =
x^\nssucc = (x \nsplus y)^\nssucc$. If $x$ and $y$ are both
non-standard, we have four cases:
\begin{align*}
&  a \nsplus a^\nssucc  = b = b^\nssucc = (a \nsplus a)^\nssucc\\
&  b \nsplus b^\nssucc  = a = a^\nssucc = (b \nsplus b)^\nssucc\\
&  b \nsplus a^\nssucc  = b = b^\nssucc = (b \nsplus y)^\nssucc\\
&  a \nsplus b^\nssucc  = a = a^\nssucc = (a \nsplus b)^\nssucc\\
\intertext{If $x$ is standard, but $y$ is non-standard, we have}
&  n \nsplus a^\nssucc  = n \nsplus a = b = b^\nssucc = (n \nsplus a)^\nssucc\\
&  n \nsplus b^\nssucc  = n \nsplus b = a = a^\nssucc = (n \nsplus b)^\nssucc
\end{align*}
So, $\Sat{L}{Q_5}$. However, $a \nsplus 0 \neq 0 \nsplus a$, so
$\Sat/{L}{\lforall[x][\lforall[y][\eq[(x+y)][(y+x)]]]}$.
\end{ex}

\begin{prob}
Expand $\Struct{L}$ of \olref[mod][mar][mdq]{ex:model-L-of-Q} to
include $\nstimes$ and $\nsless$ that interpret~$\times$ and $<$. Show
that your structure satisfies the remaining axioms of~$\Th{Q}$,
\begin{align*}
& \lforall[x][\eq[(x \times \Obj 0)][\Obj 0]] \tag{$Q_6$}\\ &
  \lforall[x][\lforall[y][\eq[(x \times y')][((x \times y) + x)]]]
  \tag{$Q_7$}\\ & \lforall[x][\lforall[y][(x < y \liff \lexists[z][\eq[(z'+x)][y]])]] \tag{$Q_8$}
\end{align*}
\end{prob}

\begin{prob}
In $\Struct{L}$ of \olref[mod][mar][mdq]{ex:model-L-of-Q}, $a^\nssucc
= a$ and $b^\nssucc = b$. Is there a model of~$\Th{Q}$ in which
$a^\nssucc = b$ and $b^\nssucc = a$?
\end{prob}

\begin{explain}
We've explicitly constructed models of~$\Th{Q}$ in which the
non-standard \olup{element}s live ``beyond'' the standard elements. In
fact, that much is required by the axioms. A non-standard
\olup{element}~$x$ cannot be ${} \nsless 0$, since $\Th{Q} \Proves
\lforall[x][\lnot x<0]$ (see \olref[inc][req][min]{lem:less-zero}).
Also, for every $n$, $\Th{Q} \Proves \lforall[x][(x < \num{n}' \lif
(\eq[x][\num{0}] \lor \eq[x][\num{1}] \lor \dots \lor
\eq[x][\num{n}]))]$ (\olref[inc][req][min]{lem:less-nsucc}), so we
can't have $a \nsless n$ for any~$n>0$.
\end{explain}


% ==== OLP file: content/model-theory/models-of-arithmetic/non-standard-models.tex ====
% Part: model-theory
% Chapter: models-of-arithmetic
% Section: non-standard-models



\olfileid{mod}{mar}{nst}
\section{Non-Standard Models}

\begin{explain}
We call \olup{structure} for $\Lang{L_A}$ standard if it is isomorphic
to~$\Struct{N}$. If \olup{structure} isn't isomorphic to~$\Struct{N}$,
it is called non-standard.
\end{explain}

\begin{defn}
A \olup{structure}~$\Struct{M}$ for $\Lang{L_A}$ is \emph{non-standard}
if it is not isomorphic to~$\Struct{N}$. The \olup{element}s $x \in
\Domain{M}$ which are equal to $\Value{\num{n}}{M}$ for some $n \in
\Nat$ are called \emph{standard numbers} (of $\Struct{M}$), and those
not, \emph{non-standard numbers}.
\end{defn}

\begin{explain}
By \olref[stm]{prop:standard-domain}, any standard \olup{structure}
for~$\Lang{L_A}$ contains only standard \olup{element}s. Consequently, a
non-standard \olup{structure} must contain at least one non-standard
element. In fact, the existence of a non-standard \olup{element}
guarantees that the \olup{structure} is non-standard.
\end{explain}

\begin{prop}
If \olup{structure}~$\Struct{M}$ for $\Lang{L_A}$ contains a
non-standard number, $\Struct{M}$ is non-standard.
\end{prop}

\begin{proof}
Suppose not, i.e., suppose $\Struct{M}$ standard but contains a
non-standard number~$x$. Let $g\colon \Nat \to \Domain{M}$ be an
isomorphism. It is easy to see (by induction on~$n$) that
$g(\Value{\num{n}}{N}) = \Value{\num{n}}{M}$. In other words, $g$ maps
standard numbers of~$\Struct{N}$ to standard numbers
of~$\Struct{M}$. If $\Struct{M}$ contains a non-standard number, $g$
cannot be \olup{surjective}, contrary to hypothesis.
\end{proof}

\begin{prob}
Recall that $\Th{Q}$ contains the axioms
\begin{align*}
& \lforall[x][\lforall[y][(\eq[x'][y'] \lif \eq[x][y])]] \tag{$Q_1$}\\
& \lforall[x][\eq/[\Obj 0][x']] \tag{$Q_2$}\\
& \lforall[x][(\eq[x][\Obj 0] \lor \lexists[y][\eq[x][y']])] \tag{$Q_3$}
\end{align*}
Give \olup{structure}s~$\Struct{M_1}$, $\Struct{M_2}$, $\Struct{M_3}$ such that
\begin{enumerate}
\item $\Sat{M_1}{Q_1}$, $\Sat{M_1}{Q_2}$, $\Sat/{M_1}{Q_3}$;
\item $\Sat{M_2}{Q_1}$, $\Sat/{M_2}{Q_2}$, $\Sat{M_2}{Q_3}$; and
\item $\Sat/{M_3}{Q_1}$, $\Sat{M_3}{Q_2}$, $\Sat{M_3}{Q_3}$;
\end{enumerate}
Obviously, you just have to specify~$\Assign{\Obj{0}}{M_i}$ and
$\Assign{\prime}{M_i}$ for each.
\end{prob}

\begin{explain}
It is easy enough to specify non-standard \olup{structure}s for
$\Lang{L_A}$. For instance, take the structure with \olup{domain}~$\Int$
and interpret all non-logical symbols as usual. Since negative numbers
are not values of $\num{n}$ for any~$n$, this structure is
non-standard. Of course, it will not be a \emph{model} of arithmetic
in the sense that it makes the same sentences true
as~$\Struct{N}$. For instance, $\lforall[x][\eq/[x'][\Obj{0}]]$ is
false.  However, we can prove that non-standard models of arithmetic
exist easily enough, using the compactness theorem.
\end{explain}

\begin{prop}
Let $\Th{TA} = \Setabs{A}{\Sat{N}{A}}$ be the theory
of~$\Struct{N}$. $\Th{TA}$ has \olup{enumerable} non-standard model.
\end{prop}

\begin{proof}
Expand $\Lang{L_A}$ by a new \olup{constant}~$c$ and consider the set of
\olup{sentence}s
\[
\Gamma = \Th{TA} \cup \{\eq/[c][\num{0}], \eq/[c][\num{1}],
\eq/[c][\num{2}], \dots\}
\]
Any model~$\Struct{M^c}$ of~$\Gamma$ would contain \olup{element}~$x =
\Assign{c}{M}$ which is non-standard, since $x \neq
\Value{\num{n}}{M}$ for all $n \in \Nat$. Also, obviously,
$\Sat{M^c}{\Th{TA}}$, since $\Th{TA} \subseteq \Gamma$. If we turn
$\Struct{M^c}$ into \olup{structure}~$\Struct{M}$ for $\Lang{L_A}$
simply by forgetting about~$c$, its domain still contains the
non-standard~$x$, and also~$\Sat{M}{\Th{TA}}$. The latter is
guaranteed since $c$ does not occur in~$\Th{TA}$. So, it suffices to
show that $\Gamma$ has a model.

We use the compactness theorem to show that~$\Gamma$ has a model. If
every finite subset of~$\Gamma$ is satisfiable, so
is~$\Gamma$. Consider any finite subset $\Gamma_0 \subseteq
\Gamma$. $\Gamma_0$ includes some \olup{sentence}s of~$\Th{TA}$ and some
of the form~$\eq/[c][\num{n}]$, but only finitely many. Suppose $k$ is
the largest number so that $\eq/[c][\num{k}] \in \Gamma_0$. Define
$\Struct{N_k}$ by expanding~$\Struct{N}$ to include the
interpretation~$\Assign{c}{N_k} = k+1$. $\Sat{N_k}{\Gamma_0}$: if $A
\in \Th{TA}$, $\Sat{N_k}{A}$ since $\Struct{N_k}$ is just
like~$\Struct{N}$ in all respects except~$c$, and $c$ does not occur
in~$A$. And $\Sat{N_k}{\eq/[c][\num{n}]}$, since $n \le k$, and
$\Value{c}{N_k} = k+1$. Thus, every finite subset of~$\Gamma$ is
satisfiable.
\end{proof}


% ==== OLP file: content/model-theory/models-of-arithmetic/standard-models.tex ====
% Part: model-theory
% Chapter: models-of-arithmetic
% Section: standard-models



\olfileid{mod}{mar}{stm}
\olsection{Standard Models of Arithmetic}

The language of arithmetic~$\Lang{L_A}$ is obviously intended to be
about numbers, specifically, about natural numbers. So, ``the''
standard model~$\Struct{N}$ is special: it is the model we want to
talk about. But in logic, we are often just interested in structural
properties, and any two \olup{structure}s that are isomorphic share
those.  So we can be a bit more liberal, and consider any
\olup{structure} that is isomorphic to~$\Struct{N}$ ``standard.''

\begin{defn}
A \olup{structure} for $\Lang{L_A}$ is \emph{standard} if it is
isomorphic to~$\Struct{N}$.
\end{defn}

\begin{prop}
\ollabel{prop:standard-domain} If \olup{structure}~$\Struct{M}$ is standard,
then its domain is the set of values of the standard numerals, i.e.,
\[
\Domain{M} = \Setabs{\Value{\num{n}}{M}}{n \in \Nat}
\]
\end{prop}

\begin{proof}
Clearly, every $\Value{\num{n}}{M} \in \Domain{M}$. We just have to
show that every $x \in \Domain{M}$ is equal to $\Value{\num{n}}{M}$
for some~$n$.  Since $\Struct{M}$ is standard, it is isomorphic
to~$\Struct{N}$. Suppose $g\colon \Nat \to \Domain{M}$ is an
isomorphism. Then $g(n) = g(\Value{\num{n}}{N}) =
\Value{\num{n}}{M}$. But for every $x \in \Domain{M}$, there is an~$n
\in \Nat$ such that $g(n) = x$, since $g$ is \olup{surjective}.
\end{proof}

\begin{explain}
If a structure~$\Struct{M}$ for $\Lang{L_A}$ is standard, the elements
of its \olup{domain} can all be named by the standard numerals $\num{0}$,
$\num{1}$, $\num{2}$, \dots, i.e., the terms $\Obj{0}$, $\Obj{0}'$,
$\Obj{0}''$, etc. Of course, this does not mean that the \olup{element}s
of $\Domain{M}$ \emph{are} the numbers, just that we can pick them out
the same way we can pick out the numbers in~$\Domain{N}$.
\end{explain}

\begin{prob}
Show that the converse of \olref[mod][mar][stm]{prop:standard-domain}
is false, i.e., give an example of \olup{structure}~$\Struct{M}$ with
$\Domain{M} = \Setabs{\Value{\num{n}}{M}}{n \in \Nat}$ that is not
isomorphic to~$\Struct{N}$.
\end{prob}

\begin{prop}
\ollabel{prop:thq-standard}
If $\Sat{M}{\Th{Q}}$, and $\Domain{M} = \Setabs{\Value{\num{n}}{M}}{n
  \in \Nat}$, then $\Struct{M}$ is standard.
\end{prop}

\begin{proof}
We have to show that $\Struct{M}$ is isomorphic
to~$\Struct{N}$. Consider the function $g\colon \Nat \to \Domain{M}$
defined by $g(n) = \Value{\num{n}}{M}$. By the hypothesis, $g$ is
\olup{surjective}.  It is also \olup{injective}: $\Th{Q} \Proves
\eq/[\num{n}][\num{m}]$ whenever $n \neq m$. Thus, since
$\Sat{M}{\Th{Q}}$, $\Sat{M}{\eq/[\num{n}][\num{m}]}$, whenever $n \neq
m$. Thus, if $n \neq m$, then $\Value{\num{n}}{M} \neq
\Value{\num{m}}{M}$, i.e., $g(n) \neq g(m)$.

We also have to verify that $g$ is an isomorphism.
\begin{enumerate}
\item We have $g(\Assign{\Obj{0}}{N}) = g(0)$ since,
  $\Assign{\Obj{0}}{N} = 0$. By definition of~$g$, $g(0) =
  \Value{\num{0}}{M}$. But $\num{0}$ is just $\Obj{0}$, and the value
  of a term which happens to be \olup{constant} is given by what the
  \olup{structure} assigns to that \olup{constant}, i.e.,
  $\Value{\Obj{0}}{M} = \Assign{\Obj{0}}{M}$. So we have
  $g(\Assign{\Obj{0}}{N}) = \Assign{\Obj{0}}{M}$ as required.
\item $g(\Assign{\prime}{N}(n)) = g(n+1)$, since $\prime$ in
  $\Struct{N}$ is the successor function on~$\Nat$. Then, $g(n+1) =
  \Value{\num{n+1}}{M}$ by definition of~$g$. But $\num{n+1}$ is the
  same term as $\num{n}'$, so $\Value{\num{n+1}}{M} =
  \Value{\num{n}'}{M}$. By the definition of the value function, this
  is $= \Assign{\prime}{M}(\Value{\num{n}}{M})$. Since
    $\Value{\num{n}}{M} = g(n)$ we get $g(\Assign{\prime}{N}(n)) =
  \Assign{\prime}{M}(g(n))$.
\item $g(\Assign{+}{N}(n,m)) = g(n+m)$, since $+$ in $\Struct{N}$ is
  the addition function on~$\Nat$. Then, $g(n+m) =
  \Value{\num{n+m}}{M}$ by definition of~$g$. But $\Th{Q} \Proves
  \num{n+m} = (\num{n} + \num{m})$, so $\Value{\num{n+m}}{M} =
  \Value{\num{n}+\num{m}}{M}$. By the definition of the value
  function, this is $=
  \Assign{+}{M}(\Value{\num{n}}{M},\Value{\num{m}}{M})$. Since
    $\Value{\num{n}}{M} = g(n)$ and $\Value{\num{m}}{M} = g(m)$, we
    get $g(\Assign{+}{N}(n, m)) = \Assign{+}{M}(g(n), g(m))$.
\item $g(\Assign{\times}{N}(n, m)) = \Assign{\times}{M}(g(n), g(m))$:
  Exercise.
\item $\tuple{n,m} \in \Assign{<}{N}$ iff $n < m$. If $n < m$, then
  $\Th{Q} \Proves \num{n} < \num{m}$, and also $\Sat{M}{\num{n} <
  \num{m}}$. Thus $\tuple{\Value{\num{n}}{M}, \Value{\num{m}}{M}} \in
  \Assign{<}{M}$, i.e., $\tuple{g(n), g(m)} \in \Assign{<}{M}$. If $n
  \not< m$, then $\Th{Q} \Proves \lnot \num{n} < \num{m}$, and
  consequently $\Sat/{M}{\num{n} < \num{m}}$. Thus, as before,
  $\tuple{g(n), g(m)} \notin \Assign{<}{M}$. Together, we get:
  $\tuple{n,m} \in \Assign{<}{N}$ iff $\tuple{g(n), g(m)} \in
  \Assign{<}{M}$.
\end{enumerate}
\end{proof}

\begin{explain}
The function~$g$ is the most obvious way of defining a mapping from
$\Nat$ to the domain of any other \olup{structure}~$\Struct{M}$ for
$\Lang{L_A}$, since every such $\Struct{M}$ contains \olup{element}s named
by $\num{0}$, $\num{1}$, $\num{2}$, etc. So it isn't surprising that
if $\Struct{M}$ makes at least some basic statements about the
$\num{n}$'s true in the same way that~$\Struct{N}$ does, and $g$ is
also bijective, then $g$ will turn into an isomorphism.  In fact, if
$\Domain{M}$ contains no \olup{element}s other than what the $\num{n}$'s
name, it's the only one.
\end{explain}

\begin{prop}
\ollabel{prop:thq-unique-iso} If $\Struct{M}$ is standard, then $g$ from
the proof of \olref{prop:thq-standard} is the only isomorphism from $\Struct{N}$ to~$\Struct{M}$.
\end{prop}

\begin{proof}
Suppose $h\colon \Nat \to \Domain{M}$ is an isomorphism between
$\Struct{N}$ and~$\Struct{M}$. We show that $g = h$ by induction
on~$n$. If $n = 0$, then $g(0) = \Assign{\Obj{0}}{M}$ by definition
of~$g$. But since $h$ is an isomorphism, $h(0) =
h(\Assign{\Obj{0}}{N}) =\Assign{\Obj{0}}{M}$, so $g(0) = h(0)$.

Now consider the case for $n+1$. We have
\begin{align*}
  g(n+1) & = \Value{\num{n+1}}{M} \text{ by definition of~$g$}\\
  & = \Value{\num{n}'}{M} \text{ since $\num{n+1}\ident \num{n}'$}\\
  & = \Assign{\prime}{M}(\Value{\num{n}}{M}) 
    \text{ by definition of $\Value{t'}{M}$}\\
  & = \Assign{\prime}{M}(g(n))  \text{ by definition of~$g$}\\
  & = \Assign{\prime}{M}(h(n)) \text{ by induction hypothesis}\\
  & = h(\Assign{\prime}{N}(n)) \text{ since $h$ is an isomorphism}\\
  & = h(n+1)
\end{align*}
\end{proof}

\begin{explain}
For any \olup{denumerable} set~$M$, there's \olup{bijection} between
$\Nat$ and $M$, so every such set~$M$ is potentially the \olup{domain} of
a standard model~$\Struct{M}$. In fact, once you pick an object $z \in
M$ and a suitable function $s$ as $\Assign{\Obj{0}}{M}$ and
$\Assign{\prime}{M}$, the interpretations of $+$, $\times$, and $<$ is
already fixed.  Only functions~$s\colon M \to M \setminus \{z\}$ that
are both \olup{injective} and \olup{surjective} are suitable in a standard
model as~$\Assign{\prime}{M}$.  The range of $s$ cannot contain~$z$,
since otherwise $\lforall[x][\eq/[\Obj 0][x']]$ would be false. That
\olup{sentence} is true in~$\Struct{N}$, and so $\Struct{M}$ also has to
make it true. The function~$s$ has to be \olup{injective}, since the
successor function~$\Assign{\prime}{N}$ in~$\Struct{N}$ is, and that
$\Assign{\prime}{N}$ is \olup{injective} is expressed by \olup{sentence}
true in~$\Struct{N}$. It has to be \olup{surjective} because otherwise
there would be some $x \in M \setminus \{z\}$ not in the domain
of~$s$, i.e., the \olup{sentence} $\lforall[x][(\eq[x][\Obj 0] \lor
\lexists[y][\eq[y'][x]])]$ would be false in~$\Struct{M}$---but it is
true in~$\Struct{N}$.
\end{explain}



\end{document}
